\documentclass[10pt,a4paper]{amsart}
\usepackage[margin=19mm]{geometry}
\usepackage{amsmath,amssymb,mathtools}
\usepackage[scr=rsfs,cal=euler]{mathalfa}
\usepackage{xfrac}
\usepackage[unicode,hidelinks,bookmarks=false]{hyperref}
\newcommand{\ie}{i.e.\ }
\newcommand{\cf}{cf.\ }
\newcommand{\resp}{respectively\ }
\newcommand{\Subsection}{\subsection}
\providecommand{\nobreakdash}{\nobreak\hbox{-}}
\setlength{\emergencystretch}{2em}
\multlinegap0pt
\usepackage{amssymb}
\usepackage[shortlabels]{enumitem}
\usepackage{tikz}
\newtheorem{proposition}{Proposition}
\newtheorem{lemma}{Lemma}
\title{Every natural number is a sum of distinct semiprime unit fractions}
\author{Shisheng Li}
\date{Source: arXiv:2606.15159v2 (CC BY 4.0)}
\begin{document}
\maketitle

\noindent\emph{Source and licence.} Every natural number is a sum of distinct semiprime unit fractions. Version arXiv:2606.15159v2 (CC BY 4.0), distributed by arXiv under the Creative Commons Attribution 4.0 International licence, http://creativecommons.org/licenses/by/4.0/, as recorded in the arXiv licence metadata for this exact version and shown on the abstract page https://arxiv.org/abs/2606.15159v2. Full licence notice: \texttt{licenses/arxiv-2606.15159-CC-BY-4.0.txt}.\par\smallskip

\noindent\textit{Editorial scope.} Counted unit: the article's lemma lem:tail-elem (source0.tex lines 831--834). The article's complete original proof of it is delivered with it (source0.tex lines 836--878, one \texttt{proof} environment of 43 lines). No mathematics is written, repaired or re-derived: the statement and the proof are the article's own lines, in the article's own notation, and no retained line is substituted or re-flowed.  Retained uncounted as the source-local context the proof needs: lines 399--405.  Nothing was omitted from any retained span, so the package carries no omission record.  No retained line contains a citation, so the wrapper carries no bibliography and no bibliography entry.  No retained line contains a figure environment, an image-inclusion command or drawing code, so no image is delivered and the package declares no asset dependency.  Editorial formatting only, in addition: an AMS \texttt{amsart} preamble that restates the article's own declarations and macros used by the retained lines, namely the environments \texttt{proposition}, \texttt{lemma} on separate counters, together with the standard packages those lines need.  The retained environments print in this document as Proposition 1, Lemma 1 in order of appearance, whereas the article prints the same items with the numbering of its own full document.  The complete native source of the article as distributed in the arXiv e-print for this exact version is bundled unchanged as \texttt{sources/202833/source0.tex}.
\begin{proposition}[explicit overlap]\label{prop:overlap}
For every $N\ge 1$ the sequence $B_N$ is strictly increasing, and for every
$N\ge 10$,
\[
B_{N+1}\le 5B_N,\qquad\text{equivalently}\qquad B_{N+1}-B_N\le 4B_N .
\]
\end{proposition}

\begin{lemma}[elementary certified tail]\label{lem:tail-elem}
For all $N\ge 300$, inequality \textup{(Tail)} holds.  Consequently
$\sigma(C)\le\tfrac16\sigma^1(N)-q$ and $Y_0(N)\le\tfrac16\sigma^1(N)-q$.
\end{lemma}

\begin{proof}
Write $g(N):=2\sqrt{U(N+1)}+1$.

\emph{Lower bound for the right side.}  $A_N$ is strictly increasing
(Proposition~\ref{prop:overlap}), so $A_N\ge A_{300}$; and
$A_{300}=\sum_{i\le 300}1/p_i=2.2909\ldots$ is a finite sum, giving
$A_N/6\ge A_{300}/6=0.3818\ldots$

\emph{Upper bound for the left side.}  The $c_0$ summands are each $\le1/p_{N-c_0}$,
so the feed tail is $\le c_0/p_{N-c_0}$.  By (Ch), $q=p_{N+1}<U(N+1)$, hence
$c_0\le 2\sqrt q+1<2\sqrt{U(N+1)}+1=g(N)$; and since $L$ is increasing and
$c_0<g(N)$, $p_{N-c_0}>L(N-c_0)\ge L(N-g(N))$.  Thus the feed tail is
$<\tau(N):=g(N)/L(N-g(N))$.  The carry is $q/P_N\le U(N+1)/2^N=:\varepsilon(N)$,
with $\varepsilon(300)<3436/2^{300}<10^{-86}$.  Hence
$\Sigma(N)<\overline\Sigma(N):=\tau(N)+\varepsilon(N)$.

\emph{Monotonicity.}  $\varepsilon$ is decreasing.  For $\tau$, with $D:=N-g(N)$,
\[
  (\log\tau)'=\frac{g'(N)}{g(N)}-\bigl(1-g'(N)\bigr)\frac{L'(D)}{L(D)}
  \ \le\ \frac{g'(N)}{g(N)}-\frac{1-g'(N)}{D},
\]
using $L'(D)/L(D)=(\log D+1)/(D\log D)\ge 1/D$ and $1-g'(N)>0$.  Since $D+g(N)=N$,
the right side is $<0$ exactly when $g'(N)\,(D+g(N))<g(N)$, i.e.\ $g'(N)<g(N)/N$.
From $g(N)=2\sqrt{U(N+1)}+1$ with $U(k)=2k\log k$ one computes
$g'(N)=2(\log(N+1)+1)/\sqrt{U(N+1)}$, which is decreasing; at $N=300$,
$g'(300)=13.41/58.62<0.23$ while $g(300)/300>118/300>0.39$, and the ratio
$g'(N)/(g(N)/N)\to\tfrac12$.  Hence $g'(N)<g(N)/N$ throughout $[300,\infty)$, so
$(\log\tau)'<0$ and $\overline\Sigma$ is decreasing there.

\emph{Base evaluation at $N=300$.}  $U(301)=602\log 301<3436$, so
$g(300)<2\sqrt{3436}+1<118.3$, $D(300)>181.7$, and
$L(D(300))>\tfrac12\cdot181.7\cdot\log 181.7>472$; whence
$\tau(300)<118.3/472<0.2507$, and adding $\varepsilon(300)<10^{-86}$ gives
$\overline\Sigma(300)<0.2507$.

\emph{Conclusion.}  For all $N\ge300$,
\[
  \Sigma(N)<\overline\Sigma(N)\le\overline\Sigma(300)<0.2507<0.3818
  =A_{300}/6\le A_N/6,
\]
which is (Tail), with margin $>0.13$.  Dividing through by $P_N$ in (O) gives
$\sigma(C)\le\tfrac16\sigma^1(N)-q$ and hence $Y_0(N)\le\tfrac16\sigma^1(N)-q$.
\end{proof}

\end{document}
